\documentclass[11pt,a4paper]{article}
% Shared preamble for RuPhO English problem statements (match T1 2026 style).
% Use: \input{latex/rupho_en_preamble.tex} after \documentclass.
\usepackage[margin=1in]{geometry}
\usepackage{amsmath,amssymb}
\usepackage{graphicx}
\usepackage{float}
\usepackage{enumitem}
\usepackage{microtype}
\usepackage{caption}
\usepackage{tabularx}
\usepackage{hyperref}

\captionsetup{font=small,labelfont=bf,skip=6pt}

\setlist[itemize]{leftmargin=1.2cm}
\setlist[enumerate]{leftmargin=1.2cm}

% One outer box per subproblem: [id | statement | points] (no inner rules)
\newcommand{\problembox}[3]{%
  \par\addvspace{0.42em}%
  \noindent
  \setlength{\fboxsep}{7.5pt}%
  \setlength{\fboxrule}{0.95pt}%
  \fbox{%
    \begin{minipage}{\dimexpr\linewidth-2\fboxsep-2\fboxrule\relax}
    \begin{tabularx}{\linewidth}{@{}>{\raggedright\arraybackslash\bfseries}p{2.1em} @{\hspace{0.9em}} >{\raggedright\arraybackslash}X @{\hspace{0.9em}} >{\raggedleft\arraybackslash\bfseries}p{2.4em}@{}}
    #1 & #2 & #3
    \end{tabularx}
    \end{minipage}%
  }%
  \par\addvspace{0.32em}%
}

\title{\textbf{Charge movement in a dipole field}}
\author{\normalsize Y13 (2013) --- English translation}
\date{}
\begin{document}
\maketitle

There is a rigidly fixed electric dipole (two opposite, but identical in module point charges $+q$ and $-q$, located at a distance $l=2~\text{cm}$ from each other). At a distance $L=2~\text{cm}$ from the dipole on a line perpendicular to its axis (Fig. 3), a ball of mass $m$, charged with a charge $+q$, is located at equal distances from the dipole charges. The ball is released without a push ($v_0=0$), and it begins to move in the electric field of the dipole.

\begin{figure}[H]
    \centering
    \includegraphics[width=0.72\textwidth]{T2_EN_assets/fig1.jpg}
    \caption{Figure 1.}
\end{figure}

\problembox{A1}{Show that the curve along which the ball will move is an arc of a circle.}{0}

\problembox{A2}{Indicate the points at which the ball will have maximum and minimum speeds.}{0}

\problembox{A3}{Find the maximum speed of the ball, taking the value $\dfrac{q^2}{4\pi \varepsilon_0 m}=\dfrac{\text{V}}{\text{m} \cdot \text{kg}}$}{0}

\vspace{1em}
\noindent\textit{Source:} \url{https://pho.rs/p/3972}

\end{document}